\chapter{داده‌کاوی و فضاهای با بعد بالا}
\label{ch:data-mining}

\section{مقدمه و تعاریف پایه}
در دنیای مدرن، سازمان‌ها، موتورهای جستجو، شبکه‌های اجتماعی و سامانه‌های علمی با حجم عظیمی از داده‌های دیجیتال روبرو هستند. گردآوری و ذخیره‌سازی داده‌ها گام نخست است، اما ارزش واقعی داده‌ها زمانی پدیدار می‌شود که بتوان از آن‌ها \textbf{دانش}\footnote{\lr{Knowledge}} و بینش عملیاتی استخراج نمود. رشته علمی و مهندسی که به این مقصود اختصاص دارد، \textbf{داده‌کاوی}\footnote{\lr{Data Mining}} نامیده می‌شود.

\begin{definition}[داده‌کاوی]
داده‌کاوی به فرآیند کشف خودکار یا نیمه‌خودکار الگوها، مدل‌ها و همبستگی‌های پنهان در مجموعه‌داده‌های عظیم اطلاق می‌شود، به‌گونه‌ای که الگوهای به‌دست‌آمده دارای ویژگی‌های زیر باشند:
\begin{enumerate}
    \item \textbf{معتبر بودن}\footnote{\lr{Valid}}: بر روی داده‌های جدید نیز با درجه اطمینان مشخصی صادق باشد (پرهیز از بیش‌برازش).
    \item \textbf{سودمند بودن}\footnote{\lr{Useful}}: بتوان بر مبنای آن تصمیمی عملیاتی یا تجاری اتخاذ کرد.
    \item \textbf{غیربدیهی و غیرمنتظره بودن}\footnote{\lr{Unexpected}}: صرفاً تکرار بدیهیات و روابط از پیش شناخته‌شده نباشد.
    \item \textbf{قابل فهم بودن برای انسان}\footnote{\lr{Understandable}}: ساختار کشف‌شده توسط تحلیل‌گران قابل تفسیر باشد.
\end{enumerate}
\end{definition}

از منظر اهداف تحلیلی، روش‌های داده‌کاوی غالباً به دو دسته بزرگ تقسیم می‌شوند:
\begin{itemize}
    \item \textbf{روش‌های توصیفی}\footnote{\lr{Descriptive Methods}}: یافتن الگوهای ساختاری که داده‌های موجود را به شیوه‌ای قابل درک برای انسان توصیف و دسته‌بندی می‌کنند. شاخص‌ترین نمونه این روش‌ها، \textbf{خوشه‌بندی}\footnote{\lr{Clustering}} و کاوش \textbf{قوانین انجمنی}\footnote{\lr{Association Rules}} است.
    \item \textbf{روش‌های پیش‌بینانه}\footnote{\lr{Predictive Methods}}: استفاده از متغیرهای معلوم برای پیش‌بینی مقادیر مجهول یا رخدادهای آینده. نمونه‌های بارز آن شامل الگوریتم‌های \textbf{طبقه‌بندی}\footnote{\lr{Classification}}، \textbf{رگرسیون}\footnote{\lr{Regression}} و \textbf{سامانه‌های توصیه‌گر}\footnote{\lr{Recommender Systems}} است.
\end{itemize}

\section{داده‌کاوی در تقاطع پایگاه‌های داده، یادگیری ماشین و نظریه الگوریتم‌ها}
داده‌کاوی رشته‌ای مجزا و در عین حال پیونددهنده میان چند حوزه کلاسیک در علوم کامپیوتر است (شکل \ref{fig:dm-intersection}):

\begin{figure}[H]
\centering
\begin{tikzpicture}[scale=0.95, every node/.style={transform shape}]
    \draw[thick, fill=blue!12, opacity=0.7] (-1.6, 0) circle (2.5cm);
    \draw[thick, fill=green!12, opacity=0.7] (1.6, 0) circle (2.5cm);
    \draw[thick, fill=orange!12, opacity=0.7] (0, -2.2) circle (2.5cm);
    
    \node at (-1.6, 0.7) [align=center] {\rl{\small\textbf{سیستم‌های پایگاه داده}}\\[3pt]\rl{\footnotesize داده‌های بسیار حجیم}\\[2pt]\rl{\footnotesize پرس‌وجوهای ساخت‌یافته}};
    
    \node at (1.6, 0.7) [align=center] {\rl{\small\textbf{یادگیری ماشین}}\\[3pt]\rl{\footnotesize مدل‌های آماری پیچیده}\\[2pt]\rl{\footnotesize استنتاج و تعمیم}};
    
    \node at (0, -2.85) [align=center] {\rl{\small\textbf{نظریه الگوریتم‌ها}}\\[3pt]\rl{\footnotesize الگوریتم‌های تصادفی}\\[2pt]\rl{\footnotesize و روش‌های تقریبی}};
    
    \node at (0, -0.85) [align=center] {\rl{\small\textbf{داده‌کاوی}}\\\rl{\small\textbf{و کلان‌داده}}};
\end{tikzpicture}
\caption{جایگاه داده‌کاوی در تقاطع پایگاه داده، یادگیری ماشین و نظریه الگوریتم‌ها}
\label{fig:dm-intersection}
\end{figure}

تفاوت فرهنگ‌های علمی در این زمینه‌ها شایان توجه است:
\begin{enumerate}
    \item \textbf{دیدگاه پایگاه داده}\footnote{\lr{Database Culture}}: در این دیدگاه، داده‌کاوی به مثابه یک پردازش تحلیلی در مقیاس بسیار وسیع \enparen{OLAP} تلقی می‌شود؛ هدف، اجرای بهینه پرس‌وجوهایی است که میلیون‌ها رکورد را غربال می‌کنند و خروجی نهایی، پاسخ صریح همان پرس‌وجو است.
    \item \textbf{دیدگاه یادگیری ماشین}\footnote{\lr{Machine Learning Culture}}: در این فرهنگ، هدف استنتاج یک مدل ریاضی و تخمین پارامترهای آن مدل است. خروجی نهایی، بردار وزن‌ها، درخت تصمیم یا ساختار شبکه عصبی است.
    \item \textbf{دیدگاه کلان‌داده \enparen{MMDS}}: در درس کلان‌داده، ما هر دو دیدگاه را ادغام می‌کنیم، اما تاکید اصلی بر \textbf{مقیاس‌پذیری}\footnote{\lr{Scalability}}، \textbf{معماری‌های محاسباتی}\footnote{\lr{Computing Architectures}} (نظیر پردازش موازی و جریانی) و \textbf{پیچیدگی محاسباتی} است.
\end{enumerate}

\section{محدودیت‌های آماری در داده‌کاوی و اصل بونفرونی}
یکی از بزرگ‌ترین خطرات در کاوش کلان‌داده‌ها این است که تحلیل‌گر با جستجو در میان انبوهی از متغیرها، الگوهایی را کشف کند که کاملاً تصادفی و ساختگی هستند، اما از نظر ظاهری بسیار مهم جلوه می‌کنند. این چالش در علم آمار تحت عنوان \textbf{اصل بونفرونی}\footnote{\lr{Bonferroni's Principle}} شناخته می‌شود.

\begin{theorem}[اصل بونفرونی]
اگر شما برای یافتن الگوهای جالب، فضایی بسیار بزرگ‌تر از آنچه که اندازه داده‌هایتان می‌تواند پشتیبانی کند را جستجو نمایید، با قطعیت ریاضی الگوهایی پیدا خواهید کرد که صرفاً حاصل تصادف هستند و هیچ رابطه علت و معلولی یا اعتبار تجربی در جهان واقعیت ندارند.
\end{theorem}

به بیان دقیق‌تر، اگر یک فرضیه وجود داشته باشد و احتمال رخداد تصادفی آن در یک آزمایش $p$ باشد، چنانچه $N$ آزمایش مستقل انجام دهیم، امید ریاضی تعداد دفعاتی که این فرضیه به طور تصادفی ارضا می‌شود برابر است با:
\begin{equation}
E = N \times p
\end{equation}
اگر $N \times p \gg 1$ باشد، مشاهده چنین الگویی نباید دلیلی بر کشف پدیده‌ای واقعی قلمداد شود، زیرا رخداد آن بر اساس تصادف صرف حتمی است.

\begin{example}[محاسبه هتل‌های مشکوک به سبک بونفرونی]
فرض کنید سازمان امنیت ملی می‌خواهد افرادی را که مظنون به اقدامات هماهنگ تروریستی هستند با ردیابی محل اقامت آنها شناسایی کند. سناریوی زیر در اسلایدهای درس و فصل اول کتاب مطرح شده است:
\begin{itemize}
    \item تعداد افراد تحت نظر: $10^9$ نفر (یک میلیارد نفر).
    \item بازه زمانی رصد: $1000$ روز.
    \item هر فرد به طور میانگین در $1\%$ روزها (یعنی از هر ۱۰۰ روز، ۱ روز) در یک هتل اقامت می‌کند. بنابراین در طول ۱۰۰۰ روز، هر فرد به طور تصادفی در ۱۰ روز به هتل می‌رود.
    \item هر هتل گنجایش دقیقاً $100$ مسافر دارد.
    \item چون در هر روز $10^9 \times 0.01 = 10^7$ نفر در هتل‌ها هستند، به تعداد $10^7 / 100 = 10^5$ هتل نیاز است.
    \item تعریف مظنون: هر زوج از افراد نامرتبط که \textbf{حداقل در دو روز متمایز} در یک هتل یکسان اقامت کرده باشند.
\end{itemize}
پرسش: اگر همه افراد کاملاً رفتاری تصادفی داشته باشند و هیچ توطئه‌ای در کار نباشد، چه تعداد «زوج مظنون» در این داده‌ها ظاهر خواهد شد؟

\textbf{حل گام‌به‌گام:}
\begin{enumerate}
    \item \textbf{تعداد کل زوج‌های ممکن}:
    \begin{equation}
    \binom{10^9}{2} = \frac{10^9 \times (10^9 - 1)}{2} \approx 5 \times 10^{17} \text{ زوج}
    \end{equation}
    
    \item \textbf{احتمال حضور دو فرد $A$ و $B$ در یک روز خاص در هتل یکسان}:
    احتمال اینکه $A$ در هتل باشد $0.01$، و احتمال اینکه $B$ در هتل باشد نیز $0.01$ است. چون $10^5$ هتل وجود دارد، احتمال اینکه هر دو در یک هتل خاص اقامت کنند برابر با $1/10^5$ است. بنابراین، احتمال اقامت همزمان در یک هتل در یک روز معین:
    \begin{equation}
    p_1 = 10^{-2} \times 10^{-2} \times 10^{-5} = 10^{-9}
    \end{equation}
    
    \item \textbf{احتمال اقامت در دو روز یکسان}:
    احتمال اینکه این دو نفر در دو روز خاص هم‌هتل شوند $p_1^2 = 10^{-18}$ است. تعداد جفت‌روزهای ممکن از میان ۱۰۰۰ روز برابر است با:
    \begin{equation}
    \binom{1000}{2} = \frac{1000 \times 999}{2} \approx 5 \times 10^5
    \end{equation}
    بنابراین، احتمال اینکه دو فرد تصادفی حداقل در دو روز در یک هتل اقامت داشته باشند:
    \begin{equation}
    p_2 = \binom{1000}{2} \times (10^{-9})^2 = (5 \times 10^5) \times 10^{-18} = 5 \times 10^{-13}
    \end{equation}
    
    \item \textbf{امید ریاضی تعداد زوج‌های مظنون کاذب}:
    تعداد زوج‌های مظنون مورد انتظار از حاصل‌ضرب تعداد کل زوج‌ها در احتمال $p_2$ حاصل می‌شود:
    \begin{equation}
    E = (5 \times 10^{17}) \times (5 \times 10^{-13}) = 250{,}000 \text{ زوج مظنون!}
    \end{equation}
\end{enumerate}
\textbf{نتیجه‌گیری تحلیلی:} در این داده‌ها، ۲۵۰ هزار جفت فرد کاملاً بی‌گناه شناسایی می‌شوند که تصادفاً دو بار در یک هتل اقامت داشته‌اند. اگر تیم امنیتی بخواهد تمامی این ۲۵۰ هزار زوج را بازجویی کند، دچار فلج عملیاتی خواهد شد. این مثال اثبات می‌کند که در کلان‌داده، یافتن همبستگی بدون پشتوانه آماری و شواهد کمکی، به نتایج فاجعه‌بار منجر می‌شود.
\end{example}

\begin{examnote}
در آزمون‌های درس کلان‌داده، مسائل مبتنی بر اصل بونفرونی بسیار پرکاربرد هستند. به خاطر بسپارید که کلید حل این مسائل در سه گام خلاصه می‌شود:
۱. محاسبه تعداد کل ترکیبات نامزدها $(N)$؛
۲. محاسبه احتمال رخداد تصادفی رویداد مورد نظر برای یک نامزد $(p)$؛
۳. محاسبه امید ریاضی $E = N \times p$. چنانچه $E \ge 1$ باشد، الگو غیرقابل استناد به عنوان ناهنجاری است.
\end{examnote}

\section{مفاهیم کمکی و ابزارهای ضروری برای کلان‌داده}

\subsection{معیار اهمیت کلمات در اسناد: TF-IDF}
در بازیابی اطلاعات و کاوش متون، لازم است کلمات کلیدی اسناد را بدون دخالت کلمات پرتکرار و بی‌ارزش (نظیر حروف اضافه و ربط) شناسایی کرد. برای این منظور از شاخص \textbf{TF-IDF}\footnote{\lr{Term Frequency - Inverse Document Frequency}} استفاده می‌شود.

فرض کنید مجموعه‌ای از $N$ سند متنی داریم. برای واژه $t$ در سند $d$:
\begin{enumerate}
    \item \textbf{فراوانی واژه \enparen{TF}}: نسبت تعداد دفعات تکرار واژه $t$ در سند $d$ $(f_{t,d})$ به بیشترین تعداد تکرار هر واژه‌ای در همان سند:
    \begin{equation}
    TF(t, d) = \frac{f_{t,d}}{\max_{w \in d} f_{w,d}}
    \end{equation}
    این نرمال‌سازی مانع از آن می‌شود که اسناد طولانی‌تر صرفاً به خاطر طولشان امتیاز بالاتری دریافت کنند.
    
    \item \textbf{معکوس فراوانی سند \enparen{IDF}}: اگر کلمه $t$ در $n_t$ سند از کل $N$ سند ظاهر شده باشد، شاخص IDF آن به صورت زیر تعریف می‌شود:
    \begin{equation}
    IDF(t) = \log_2 \left( \frac{N}{n_t} \right)
    \end{equation}
    اگر کلمه‌ای در تمام اسناد تکرار شود $(n_t = N)$، آنگاه $\log_2(1) = 0$ شده و وزن آن کلمه صفر می‌گردد. در مقابل، کلماتی که فقط در چند سند خاص ظاهر شوند، IDF بالایی دریافت می‌کنند.
    
    \item \textbf{شاخص ترکیبی TF-IDF}:
    \begin{equation}
    \text{TF-IDF}(t, d) = TF(t, d) \times IDF(t)
    \end{equation}
\end{enumerate}

\subsection{توابع درهم‌سازی و رفتار آن‌ها}
توابع درهم‌سازی\footnote{\lr{Hash Functions}} ابزار بنیادین در تمام الگوریتم‌های کلان‌داده هستند. یک تابع درهم‌سازی $h(x)$ کلیدی دلخواه را به مقداری صحیح در بازه $[0, B-1]$ نگاشت می‌کند. در کلان‌داده، دو ویژگی حیاتی زیر مورد انتظار است:
\begin{itemize}
    \item \textbf{توزیع یکنواخت}\footnote{\lr{Uniform Distribution}}: کلیدها به طور مساوی میان $B$ سبد پخش شوند تا از تجمع بار جلوگیری گردد.
    \item \textbf{شبه‌تصادفی بودن}\footnote{\lr{Pseudorandomness}}: حتی اگر ورودی‌ها دارای الگوهای منظم باشند، مقادیر درهم‌سازی حاصل رفتاری مشابه متغیرهای تصادفی مستقل نشان دهند.
\end{itemize}

\subsection{قوانین توانی و مفهوم دنباله طولانی}
در بسیاری از پدیده‌های وب و کلان‌داده (مانند توزیع درجات در گراف وب، تعداد خریدهای کالاها در آمازون، یا تعداد بازدید ویدئوها در یوتیوب)، داده‌ها از توزیع نرمال گاوسی پیروی نمی‌کنند، بلکه پیرو \textbf{قانون توانی}\footnote{\lr{Power Law}} هستند:
\begin{equation}
y = c \cdot x^k \quad (k < 0)
\end{equation}
که در مقیاس لگاریتمی به یک خط راست تبدیل می‌شود:
\begin{equation}
\log y = \log c + k \log x
\end{equation}
این توزیع بیانگر پدیده \textbf{دنباله طولانی}\footnote{\lr{Long Tail}} است؛ یعنی تعداد بسیار کمی از عناصر دارای فراوانی نجومی هستند (اقلیت پرطرفدار)، اما مجموع حجم عظیمی از داده‌ها در دنباله‌ای از عناصر کم‌تکرار قرار دارد.

\textbf{قانون زیپف و توزیع پارتو}:
یکی از شناخته‌شده‌ترین مصادیق قوانین توانی، \textbf{قانون زیپف}\footnote{\lr{Zipf's Law}} است که بیان می‌دارد فراوانی عنصر با رتبه $r$ ام در میان اقلام به صورت زیر رفتار می‌کند:
\begin{equation}
f(r) \propto \frac{1}{r^\alpha} \quad (\alpha \approx 1)
\end{equation}
به عنوان مثال، پربسامدترین واژه در زبان انگلیسی («the») تقریباً دو برابر دومین واژه («of») و سه برابر سومین واژه («and») تکرار می‌شود. قاعده مشهور \textbf{پارتو ۲۰-۸۰} (که در آن ۸۰ درصد خریدهای یک فروشگاه تنها از ۲۰ درصد کالاهای پرفروش ناشی می‌شود) حالت خاصی از همین توزیع‌های توانی با نماهای شیب‌دار است. در سامانه‌های کلان‌داده، پردازش کارآمد داده‌های دنباله طولانی نیازمند الگوریتم‌های متمایز از پردازش اقلام ابرپرطرفدار است.

\begin{summarybox}[جمع‌بندی فصل اول]
\begin{itemize}
    \item داده‌کاوی به معنای کشف الگوهای معتبر، جدید، سودمند و قابل فهم در حجم انبوه داده‌ها است.
    \item اصل بونفرونی هشدار می‌دهد که جستجوی کورکورانه در میان تعداد حالات نجومی، قطعاً به کشف همبستگی‌های تصادفی و کاذب می‌انجامد.
    \item شاخص TF-IDF ابزاری اساسی برای ارزیابی اهمیت واژگان در اسناد است.
    \item توابع درهم‌سازی و قوانین توانی در سراسر فصول بعدی برای طراحی الگوریتم‌های مقیاس‌پذیر به کار گرفته خواهند شد.
\end{itemize}
\end{summarybox}
